% =============================================================================
% БИЛЕТ 9
% =============================================================================

\ticket{9}{}

\question{Предел последовательности точек в $n$-мерном пространстве. Связь между сходимостью покоординатно и по норме. Теорема Больцано-Вейерштрасса}

\begin{defn}[Предел по множеству по Коши]
  Пусть $E \subset \R^n$, $x^0$ --- предельная точка $E$, будем говорить, что $C \in \R$ является пределом отображения $f$ по совокупности переменных и записывать 
    $$\lim_{x \to x^0} f(x) = C \quad \text{или} \quad \lim_{\substack{x_1 \to x^0_1 \\ \vdots \\ x_n \to x^0_n}} f(x_1,\ldots,x_n) = C$$
    если $\forall \varepsilon > 0 \exists \delta > 0 : \forall x \in E \subset B(x^0, \delta) \setminus \{x^0\}  |f(x) - C| < \varepsilon$.
\end{defn}

\begin{defn}[Предел по множеству по Гейне]
  Пусть $E \subset \R^n$, $x^0$ --- предельная точка $E$, будем говорить, что $C \in \R$ является пределом отображения $f$ по совокупности переменных и записывать 
    $$\lim_{x \to x^0} f(x) = C \quad \text{или} \quad \lim_{\substack{x_1 \to x^0_1 \\ \vdots \\ x_n \to x^0_n}} f(x_1,\ldots,x_n) = C$$
    если для всякой последовательности $\{x^{(m)}\}$, такой что $x^{(m)} \in E$ и $\lim_{m \to +\infty} x^{(m)} = x^0$, выполняется $\lim_{m \to +\infty} f(x^{(m)}) = C$.
\end{defn}

\begin{theorem}[Эквивалентность определений предела по Коши и по Гейне]
  Пусть $E \subset \R^n$, $x^0$ --- предельная точка $E$, $f : E \to \R$. Тогда следующие условия эквивалентны:
  \begin{enum}
    \item $\lim_{x \to x^0} f(x) = C$ по Коши;
    \item $\lim_{x \to x^0} f(x) = C$ по Гейне.
  \end{enum}
  \begin{proof}
    \textbf{$1 \Rightarrow 2$.} Пусть $\{x^{(m)}\}$ --- произвольная последовательность в $E$, сходящаяся к $x^0$. Зафиксируем $\varepsilon > 0$. По определению (1) найдётся $\delta > 0$, такое что для всякого $x \in E$ с $0 < |x - x^0| < \delta$ выполняется неравенство $|f(x) - C| < \varepsilon$. Поскольку $\lim_{m \to +\infty} x^{(m)} = x^0$, найдётся номер $M$, такой что для всех $m > M$ выполняется неравенство $|x^{(m)} - x^0| < \delta$. Поэтому для всех $m > M$ имеем
    $$|f(x^{(m)}) - C| < \varepsilon,$$
    что означает $\lim_{m \to +\infty} f(x^{(m)}) = C$.

    \textbf{$2 \Rightarrow 1$.} Предположим противное, то есть что $\lim_{x \to x^0} f(x) = C$ по Коши не выполняется. Тогда существует $\varepsilon_0 > 0$, такое что для всякого $\delta > 0$ найдётся точка $x_\delta \in E$ с $0 < |x_\delta - x^0| < \delta$ и $|f(x_\delta) - C| \geqslant \varepsilon_0$. Выбирая $\delta = 1/m$, получаем последовательность $\{x^{(m)}\}$ \- противоречие, так как $x^{(m)} \to x^0$ и $f(x^{(m)})$ не сходится к $C$. \qedhere
  \end{proof}
\end{theorem}

\begin{defn}[предел по вектору]
Зафиксируем произвольный вектор $l = (l_1, \ldots, l_n) \in \R^n$ и будем говорить, что $C \in \R$ является пределом функции $f : E \to \R$ по вектору $l$ при $x \to x^0$, если $\lim_{t \to +0} f(x^0 + tl) = C$. Мы при помощи вектора как бы по лучу приближаемся к точке $x^0$ и смотрим, существует ли предел функции при таком приближении. Если предел по вектору $l$ существует, то его будем обозначать как $\lim_{t \to +0} f(x^0 + tl)$
\end{defn}
\begin{lemma}
  Если $\lim_{x \to x^0} f(x) = C$, то для всякого вектора $l$ предел по вектору $l$ при $x \to x^0$ существует и равен $C$. Но обратное не верно. Когда мы ищем предел по направлению, мы должны проверить все направления, а не только по вектору $l$. Поэтому если предел по вектору $l$ существует, то это не гарантирует существование предела.
  \begin{proof}
    Так как по определению предела по Коши мы в какой то момент будем около точки $x^0$ в шаре, то как только мы достигнем $\varepsilon$ равной длине вектора $l$, мы будем в этом шаре, и так как предел по Коши существует, то мы будем в $\varepsilon$ окрестности $C$. Поэтому предел по вектору $l$ будет равен $C$.
    
    Но обратное не верно, так как если мы будем приближаться к точке $x^0$ по вектору $l$, то мы можем находиться в $\varepsilon$ окрестности $C$, но если мы будем приближаться по другому вектору, то мы можем находиться в $\varepsilon$ окрестности другого числа, отличного от $C$.
  \end{proof}
\end{lemma}
\begin{example}
$f = \frac{x^2y}{x^4+y^2}$, $x^0 = (0, 0)$
можно заметить что для всяких линейных векторов $l = (l_1, l_2)$ при $t \to +0$ мы будем находиться в $\varepsilon$ окрестности $0$, так как $f(x^0 + tl) = \frac{t^4l_1^2l_2}{t^8l_1^4+t^4l_2^2} = \frac{l_1^2l_2}{t^{-4}l_1^4+l_2^2}$, и так как $t^{-4} \to +\infty$, то $\lim_{t \to +0} f(x^0 + tl) = 0$. Но если мы будем приближаться по кривой $y = x^2$, то мы будем находиться в $\varepsilon$ окрестности $\frac{1}{2}$, так как $f(x, x^2) = \frac{x^6}{2x^8} = \frac{1}{2x^2}$, и $\lim_{x \to +0} f(x, x^2) = +\infty$. Поэтому предел по Коши не существует.
\end{example}
\begin{defn}[Метод полярных координат]
  Хотим свести исследование предела функции $f : E \to \R$ при $x \to x^0$ к исследованию предела функции одной переменной. Для этого мы можем ввести полярные координаты. Пусть $x^0 = (x^0_1, x^0_2)$, тогда для всякой точки $x = (x_1, x_2)$, отличной от $x^0$, можно ввести координаты $r$ и $\varphi$ по формуле $$x_1 = x^0_1 + r\cos\varphi, \quad x_2 = x^0_2 + r\sin\varphi.$$ тут проще считать, так как $r$ --- это расстояние от точки $x$ до точки $x^0$, а $\varphi$ --- это угол между вектором $x - x^0$ и положительным направлением оси $Ox_1$. Если мы хотим исследовать предел функции $f$ при $x \to x^0$, то мы можем исследовать предел функции $g(r, \varphi) = f(x^0_1 + r\cos\varphi, x^0_2 + r\sin\varphi)$ при $r \to +0$. Если предел функции $g$ при $r \to +0$ существует и не зависит от $\varphi$, то предел функции $f$ при $x \to x^0$ существует и равен этому пределу.
\end{defn}

\begin{proof}
  докажем, что если есть предел функции $g$ при $r \to +0$, то предел функции $f$ при $x \to x^0$ существует и равен этому пределу. Зафиксируем $\varepsilon > 0$. По определению предела функции $g$ найдётся $\delta > 0$, такое что для всякого $r$ с $0 < r < \delta$ выполняется неравенство $|g(r, \varphi) - C| < \varepsilon$. Теперь возьмём произвольную точку $x = (x_1, x_2)$ с $0 < |x - x^0| < \delta$. Тогда $r = |x - x^0| < \delta$, и мы можем ввести полярные координаты для точки $x$. Поэтому
$$|f(x) - C| = |g(r, \varphi) - C| < \varepsilon,$$
что означает $\lim_{x \to x^0} f(x) = C$. \qedhere

а теперь обратно, пусть у нас есть мажоранта $M(\varphi)$, такая что $|g(r, \varphi) - C| \leqslant M(\varphi)$ для всякого $r$ с $0 < r < \delta$. Если $\lim_{r \to +0} M(\varphi) = 0$, то для всякого $\varepsilon > 0$ найдётся $\delta' > 0$, такое что для всякого $r$ с $0 < r < \delta'$ выполняется неравенство $M(\varphi) < \varepsilon$. Поэтому для всякого $r$ с $0 < r < \min\{\delta, \delta'\}$ выполняется неравенство $|g(r, \varphi) - C| < \varepsilon$, что означает $\lim_{r \to +0} g(r, \varphi) = C$. \qedhere
\end{proof}

\begin{theorem}[Связь поокоординатной сходимости и сходимости по норме]
  Пусть $E \subset \R^n$, $x^0$ --- предельная точка $E$, $f : E \to \R$. Тогда следующие условия эквивалентны:
  \begin{enum}
    \item $\lim_{x \to x^0} f(x) = C$;
    \item для всякого $k \in \{1, \ldots, n\}$ существует предел $\lim_{x_k \to x^0_k} f(x^0_1, \ldots, x^0_{k-1}, x_k, x^0_{k+1}, \ldots, x^0_n) = C$.
  \end{enum}
  \begin{proof}
    \textbf{$1 \Rightarrow 2$.} Если $\lim_{x \to x^0} f(x) = C$, то для всякого $\varepsilon > 0$ найдётся $\delta > 0$, такое что для всякого $x$ с $0 < |x - x^0| < \delta$ выполняется неравенство $|f(x) - C| < \varepsilon$. В частности, для всякого $k$ и всякого $x_k$ с $0 < |x_k - x^0_k| < \delta$ выполняется неравенство $$|f(x^0_1, \ldots, x^0_{k-1}, x_k, x^0_{k+1}, \ldots, x^0_n) - C| < \varepsilon,$$ что означает $\lim_{x_k \to x^0_k} f(x^0_1, \ldots, x^0_{k-1}, x_k, x^0_{k+1}, \ldots, x^0_n) = C$. 
    
    \textbf{$2 \Rightarrow 1$.} Пусть у нас есть предел по каждой координате, то есть для всякого $k$ и всякого $\varepsilon > 0$ найдётся $\delta_k > 0$, такое что для всякого $x_k$ с $0 < |x_k - x^0_k| < \delta_k$ выполняется неравенство $$|f(x^0_1, \ldots, x^0_{k-1}, x_k, x^0_{k+1}, \ldots, x^0_n) - C| < \varepsilon.$$ Возьмем такую точку, чтобы она была настолько близко к $x^0$, что для всякого $k$ выполняется неравенство $|x_k - x^0_k| < \delta_k$. Тогда для всякого $k$ выполняется неравенство $$|f(x^0_1, \ldots, x^0_{k-1}, x_k, x^0_{k+1}, \ldots, x^0_n) - C| < \varepsilon,$$ что означает $\lim_{x \to x^0} f(x) = C$. \qedhere
    \end{proof}
\end{theorem}

\begin{theorem}[Больцано-Вейерштрасса]
  Пусть $E \subset \R^n$ --- ограниченная последовательность, тогда можно выделить подпоследовательность, которая сходится.
  \begin{proof}
    Доказательство по индукции по размерности. Для $n = 1$ это классическая теорема Больцано-Вейерштрасса. Предположим, что теорема верна для всех $k < n$. Пусть $E \subset \R^n$ --- ограниченная последовательность. Тогда проекция $E$ на первую координату --- это ограниченная последовательность в $\R$, поэтому по теореме Больцано-Вейерштрасса найдётся подпоследовательность $E_1 \subset E$, такая что проекция $E_1$ на первую координату сходится. Теперь проекция $E_1$ на вторую координату --- это ограниченная последовательность в $\R$, поэтому по теореме Больцано-Вейерштрасса найдётся подпоследовательность $E_2 \subset E_1$, такая что проекция $E_2$ на вторую координату сходится. Продолжая этот процесс, мы получаем подпоследовательность $E_n \subset E_{n-1} \subset \ldots \subset E_1 \subset E$, такая что проекция $E_n$ на каждую координату сходится. Поэтому $E_n$ --- это подпоследовательность, которая сходится. \qedhere
  \end{proof}
\end{theorem}
\question{Бесконечная дифференцируемость суммы степенного ряда с действительными членами в интервале сходимости. Единственность представления функции степенным рядом}

% Ответ...

